\section{Results I: Does a Small Candidate Budget Suffice?}
\label{sec:results1}

\begin{figure}[t]
\centering
\includegraphics[width=\linewidth]{fig_scatter_pilot.png}
\caption{Qualitative pilot on the bounce environment with $K{=}4$ branches and fixed appearance (pilot recipe, $m{=}4$). Each column is one context whose ball has just hit a wall: grey, the three context frames; green circles, the four possible next positions; red, 16 samples decoded by the position readout. Regression (top) returns one point near the mean, flow matching with 20 or 2 Euler steps scatters between the modes, and \nom{} places its samples on the modes.}
\label{fig:scatter}
\end{figure}

We first ask whether a one-step \nom{} world model with a small budget covers enumerated futures at all, and how it compares with multi-step and one-step alternatives at matched parameters. Figure~\ref{fig:scatter} shows the phenomenon qualitatively; Tables~\ref{tab:a2}--\ref{tab:b} quantify it on the three benchmarks.

\subsection{Non-repeating contexts (A2-div)}
\label{sec:a2}
Table~\ref{tab:a2} is the main comparison on the bounce environment with $K{=}4$ branches and randomized appearance, where every test context is new (its nearest training context is $1.9\times$ farther than the typical nearest neighbour within the training set). \imlewm{} with $m{=}16$ covers $0.93$ of the enumerated modes against an ideal-sampler bound of $0.99$, with precision $0.82$: its samples lie on the branches. FM-50 reaches $0.79$ coverage but only $0.51$ precision, because its samples scatter between the branches (Figure~\ref{fig:scatter}); FM-2, iMF and Shortcut are below it on both. The energy-score generator has the \emph{lowest} coverage of all stochastic methods ($0.36$ at $m{=}8$, $0.28$ at $m{=}16$) and, like regression, has memorized the training contexts (index $0.87$--$0.97$ against chance $0.25$). \imlewm{} with $m{=}4$ sits in between ($0.46$, index $0.75$). Section~\ref{sec:results2} is about this column. The price of $m{=}16$ is a larger error on transitions that are not stochastic ($1.55$ vs.\ $1.14$\,px for $m{=}4$, $1.09$ for regression) and $1.7\times$ the training time of $m{=}4$; inference cost is independent of $m$, $3.3$\,ms per step against $107$\,ms for FM-50.

\begin{table}[t]
\caption{Bounce with $K{=}4$ branches and randomized appearance (A2-div), 8M-parameter models, $12$k steps, three seeds (mean $\pm$ std; the energy score with $m{=}16$ uses two seeds). Coverage and precision are computed from 16 samples per held-out context decoded to the ball position; the ideal 16-sample sampler reaches coverage $0.99$. kNN recall is the feature-space recall of \citet{kynkaanniemi2019improved}. Det.\ error is the readout error on deterministic transitions. The memorization index is the realized-mode share on training contexts (chance $0.25$). Latency is per predicted step at batch size 1 on a shared RTX A5000.}
\label{tab:a2}
\begin{center}
\small
\resizebox{\linewidth}{!}{%
\begin{tabular}{lcccccccc}
\toprule
Method & NFE & Coverage $\uparrow$ & Precision $\uparrow$ & kNN recall $\uparrow$ & Det.\ err.\ (px) $\downarrow$ & Memo.\ index & ms/step & Train (h) \\
\midrule
Regression & 1 & \ms{0.15}{0.01} & \ms{0.61}{0.02} & 0.00 & \ms{1.09}{0.03} & \ms{0.99}{0.01} & 4.4 & 0.3 \\
FM-50 (FlowWM recipe) & 50 & \ms{0.79}{0.01} & \ms{0.51}{0.00} & \ms{0.91}{0.00} & \ms{2.24}{0.02} & \ms{0.29}{0.01} & 107 & 0.4 \\
FM-2 (same model) & 2 & \ms{0.60}{0.00} & \ms{0.46}{0.01} & \ms{0.56}{0.01} & \ms{1.72}{0.01} & \ms{0.29}{0.01} & 6.5 & 0.4 \\
iMF ($3\times$ steps) & 1 & \ms{0.56}{0.01} & \ms{0.23}{0.02} & \ms{0.92}{0.03} & \ms{14.5}{0.7} & \ms{0.25}{0.01} & 3.4 & 3.2 \\
Shortcut ($2\times$ steps) & 1 & \ms{0.66}{0.00} & \ms{0.42}{0.01} & \ms{0.77}{0.01} & \ms{3.60}{0.07} & \ms{0.27}{0.01} & 4.3 & 0.7 \\
Energy score, $m{=}8$ & 1 & \ms{0.36}{0.02} & \ms{0.61}{0.02} & \ms{0.77}{0.06} & \ms{1.06}{0.04} & \ms{0.87}{0.06} & 3.1 & 1.2 \\
Energy score, $m{=}16$ & 1 & \ms{0.28}{0.01} & \ms{0.61}{0.01} & \ms{0.60}{0.02} & \ms{1.06}{0.03} & \ms{0.97}{0.01} & 3.3 & 2.2 \\
\imlewm{}, $m{=}4$ & 1 & \ms{0.46}{0.01} & \ms{0.68}{0.01} & \ms{0.65}{0.01} & \ms{1.14}{0.07} & \ms{0.75}{0.01} & 3.7 & 0.6 \\
\imlewm{}, $m{=}16$ & 1 & \best{\ms{0.93}{0.01}} & \best{\ms{0.82}{0.01}} & \best{\ms{0.94}{0.01}} & \ms{1.55}{0.05} & \ms{0.29}{0.00} & 3.3 & 1.1 \\
\bottomrule
\end{tabular}}
\end{center}
\end{table}

\subsection{Bouncing Shapes under FlowWM's protocol and readout (A1)}
\label{sec:a1}
Bouncing Shapes \citep{porcher2026flowwm} is the one public benchmark with enumerable futures. We use its generator unchanged, its DINOv3 feature definition and its precision/recall error, but one eighth of its training clips and much smaller models, so we report comparisons within our pipeline and quote the paper's numbers only as a reference. Table~\ref{tab:a1} covers two prediction protocols, autoregressive one-step rollouts (our setting) and the joint prediction of 16 future frames from 16 context frames (FlowWM's setting), and two readouts, our soft-argmax probe and FlowWM's heatmap readout, which we ported. In every one of the four combinations \imlewm{} has the lowest F1 error, and the gaps are at least ten times the seed standard deviation: $6.6$ vs.\ $12.3$\,px for FM-50 in the autoregressive protocol, $9.0$ vs.\ $12.6$\,px with FlowWM's own readout in its own protocol. Because contexts repeat in this benchmark (one initial condition is seen with several futures), there is no memorization and $m{=}4$ is enough; Section~\ref{sec:results2} returns to this difference. The readout matters: the deterministic joint model moves from $5.1$ to $22.5$\,px when the readout changes, which explains most of the gap between our deterministic baseline and the $17.1$\,px reported by \citet{porcher2026flowwm}, while autoregressive models are almost unaffected because they produce a clean feature map at every step (Section~\ref{sec:manifold}). Training four times longer improves every method and preserves the ranking (one seed: \imlewm{} $m{=}4$ $4.1$, FM-2 $7.1$, FM-50 $9.0$\,px).

\begin{table}[t]
\caption{Bouncing Shapes F1 error in px (lower is better; the frame is $256$\,px), three seeds. Autoregressive: 3-frame context, 16 one-step predictions, 8M models. Joint: 16 frames predicted at once from 16 context frames, 27--29M models, 12k steps at batch 64. The last-frame-copy baseline scores $21.7$ in every column. The row from \citet{porcher2026flowwm} uses eight times our training data and a different model size and is a reference only. ``--'' marks combinations that were not run.}
\label{tab:a1}
\begin{center}
\small
\begin{tabular}{lccccc}
\toprule
 & & \multicolumn{2}{c}{Autoregressive} & \multicolumn{2}{c}{Joint 16$\to$16} \\
\cmidrule(lr){3-4}\cmidrule(lr){5-6}
Method & NFE & probe & heatmap & probe & heatmap \\
\midrule
Regression & 1 & \ms{8.09}{0.15} & \ms{8.41}{0.15} & \ms{5.14}{0.01} & \ms{22.48}{0.36} \\
FM-50 & 50 & \ms{12.28}{0.13} & \ms{14.11}{0.20} & \ms{3.86}{0.02} & \ms{12.59}{0.14} \\
FM-2 & 2 & \ms{9.95}{0.19} & \ms{12.49}{0.38} & \ms{3.83}{0.02} & \ms{14.80}{0.17} \\
iMF & 1 & \ms{16.83}{0.24} & -- & -- & -- \\
Shortcut & 1 & \ms{12.07}{0.29} & -- & -- & -- \\
Energy score, $m{=}8$ (2 seeds) & 1 & \ms{8.94}{0.20} & \ms{33.77}{1.34} & -- & -- \\
\imlewm{}, $m{=}4$ & 1 & \best{\ms{6.61}{0.18}} & \best{\ms{6.83}{0.16}} & \best{\ms{3.30}{0.02}} & \ms{11.73}{0.31} \\
\imlewm{}, $m{=}16$ & 1 & \ms{7.21}{0.10} & \ms{7.64}{0.18} & \ms{3.32}{0.07} & \best{\ms{9.04}{0.06}} \\
\midrule
FlowWM paper: FlowWM / deterministic & 50 / 1 & -- & -- & -- & 4.31 / 17.1 \\
\bottomrule
\end{tabular}
\end{center}
\end{table}

\subsection{Stochastic Push-T multi-future test (B)}
\label{sec:b}
On the Push-T variant with a hidden center of gravity, the 21M \imlewm{} with $m{=}16$ has the lowest recall error at both horizons and a precision error $6.6$\,px below FM-50 at $h{=}16$ (Table~\ref{tab:b}). We quantify uncertainty with a hierarchical bootstrap that resamples seeds and then starting states (Appendix~\ref{app:stats}). The precision advantage over FM-50, FM-2 and Shortcut is significant under that procedure ($h{=}16$: $-6.6$ $[-8.5, -4.7]$, $-5.6$ $[-7.7, -3.2]$ and $-13.2$ $[-16.3, -10.2]$\,px). The recall advantage is consistent across the 128 starting states but, with three seeds, its interval includes zero against the flow baselines ($-4.2$ $[-8.8, +1.1]$ vs.\ FM-50) and excludes it only against Shortcut ($-8.2$ $[-11.3, -4.0]$). The evidence on this benchmark is therefore an advantage in precision and a non-inferiority in recall. Two facts limit what any model can show here: all models mispredict the first contact by about $12$\,px (the true block displacement is $16$\,px), and this error does not shrink with three times more training, so differences between objectives are partly masked by a shared dynamics error; and \imlewm{} with $m{=}4$ shows the collapse pattern of Section~\ref{sec:results2}, the best precision with the worst recall among stochastic models and a rollout spread of $9$\,px against $22$\,px for $m{=}16$.

\begin{table}[t]
\caption{Stochastic Push-T multi-future test set (128 pre-contact starting states $\times$ 64 futures $\times$ 16 steps): precision and recall error over the eight T vertices in px (frame $512$\,px), 21M models, three seeds (iMF two). GT spread is the mean pairwise distance among the 64 true futures. Mode coverage is not reported: at these rollout errors no model resolves the four centre-of-gravity modes.}
\label{tab:b}
\begin{center}
\small
\begin{tabular}{lccccc}
\toprule
 & & \multicolumn{2}{c}{$h{=}8$} & \multicolumn{2}{c}{$h{=}16$} \\
\cmidrule(lr){3-4}\cmidrule(lr){5-6}
Method & NFE & precision & recall & precision & recall \\
\midrule
GT spread (reference) & & \multicolumn{2}{c}{30.8} & \multicolumn{2}{c}{48.3} \\
Regression & 1 & \ms{24.8}{1.1} & \ms{39.6}{3.4} & \ms{40.2}{3.0} & \ms{66.4}{8.2} \\
FM-50 & 50 & \ms{27.1}{1.6} & \ms{24.3}{1.8} & \ms{40.7}{1.1} & \ms{42.0}{3.6} \\
FM-2 & 2 & \ms{24.2}{1.7} & \ms{23.3}{1.7} & \ms{39.7}{1.4} & \ms{41.5}{3.8} \\
iMF & 1 & \ms{37.2}{2.6} & \ms{33.5}{3.2} & \ms{55.1}{3.7} & \ms{52.4}{3.4} \\
Shortcut & 1 & \ms{28.0}{1.0} & \ms{24.2}{0.4} & \ms{47.3}{1.9} & \ms{45.9}{1.1} \\
\imlewm{}, $m{=}4$ & 1 & \best{\ms{18.7}{0.4}} & \ms{29.2}{3.4} & \best{\ms{33.8}{3.2}} & \ms{49.1}{3.5} \\
\imlewm{}, $m{=}16$ & 1 & \ms{21.8}{1.8} & \best{\ms{22.1}{2.3}} & \ms{34.1}{0.7} & \best{\ms{37.8}{2.3}} \\
\bottomrule
\end{tabular}
\end{center}
\end{table}

\subsection{Baseline fairness and latency}
\label{sec:fairness}
Four checks guard the comparison. (i) iMF is not under-trained: a 72k-step run on A2-div saturates after 24k steps at coverage $0.58$ and precision $0.26$. (ii) Flow matching is not under-trained: with four times the steps on A1 the ranking is unchanged (Section~\ref{sec:a1}); its recipe is FlowWM's, which on A2-div raised the coverage of our initial flow baseline from $0.58$ to $0.79$. (iii) Shortcut, a stronger one-step flow generator than iMF on every benchmark, is trained to the same wall-clock budget as \imlewm{} with $m{=}16$. (iv) All models within a benchmark are within $\pm4\%$ in parameters. The remaining caveat is scale: our joint FM-50 reaches $12.6$\,px under FlowWM's readout where the paper reports $4.31$ with eight times the data and a larger model, so the claim is confined to matched small-scale compute. Latency at batch size 1 is $3$--$4$\,ms per step for every one-step model and $107$\,ms for FM-50; in the planner of Section~\ref{sec:planning}, where each CEM iteration imagines $64 \times m_p$ futures, this becomes a $50\times$ difference.
